L'exercice est constitué de deux parties indépendantes.

\medskip

\noindent {\Large \textbf{Partie 1}}

\medskip

On considère l'équation différentielle :
\[(E): y' + y = \e^{-x}\]

\begin{enumerate}
\item Soit $u$ la fonction définie sur $\R$ par $u(x) = x \e^{-x}$.

Vérifier que la fonction $u$ est une solution de l'équation différentielle $(E)$.
\item On considère l'équation différentielle $\left(E'\right): y' + y = 0$.

Résoudre l'équation différentielle $\left(E'\right)$ sur $\R$.

\item En déduire toutes les solutions de l'équation différentielle $(E)$ sur $\R$.

\item Déterminer l'unique solution $g$ de l'équation différentielle $(E)$ telle que $g(0) = 2$.
\end{enumerate}


\medskip

\noindent {\Large \textbf{Partie 2}}

\medskip

Dans cette partie, $k$ est un nombre réel fixé que l'on cherche à déterminer.

On considère la fonction $f_{k}$ définie sur $\R$ par :
\[f_{k}(x) = (x + k) \e^{-x}.\]

Soit $h$ la fonction définie sur $\R$ par
\[h(x) = \e^{-x}.\]

On note $C_{k}$ la courbe représentative de la fonction $f_{k}$ dans un repère orthogonal et $C$ la courbe représentative de la fonction $h$.

On a représenté sur le graphique en fin d'exercice les courbes $C_{k}$ et $C$ sans indiquer les unités sur les axes ni le nom des courbes.

	\begin{enumerate}
		\item Sur le graphique en fin d'exercice à rendre avec la copie, l'une des courbes est en traits pointillés, l'autre est en trait plein. Laquelle est la courbe $C$ ?

		\item En expliquant la démarche utilisée, déterminer la valeur du nombre réel $k$ et placer sur le graphique ci-dessous avec la copie l'unité sur chacun des axes du graphique.
	\end{enumerate}

\begin{center}
	\begin{tikzpicture}
		\begin{axis}[x = {20mm},y = {20mm}, xmin = -2.5, xmax = 4.5, ymin = -0.5, ymax = 3,
			xtick = \empty, ytick = \empty,
			grid = major, axis lines = center]
			\addplot [line width = 1.2pt,color = red,smooth,samples = 100,domain = -2.5:4.5 ]{exp(-x)};
			\addplot [line width = 1.2pt,color = teal,dashed,smooth,samples = 100,domain = -2.5:4.5 ]{(x + 2)*exp(-x)};
		\end{axis}
	\end{tikzpicture}
\end{center}
